% Distribución íntegra por residencia del propietario cosmológico.
\section{Raccordement géométrique entre contraction anisotrope, rebond torsionnel et région cosmologique descendante}

Le Big Bang peut être compris comme une distorsion d'intensité maximale entre état,
rang et lecteur. Ce n'est pas une explosion survenue dans un espace antérieur. C'est la
limite à laquelle la projection définissant l'extension, la distance et la durée commence à
être valide. Demander ce qui s'est produit « avant » en utilisant la même horloge peut n'avoir aucun
sens parce que l'horloge appartient au régime qui émerge. HMT conserve cependant
une généalogie antérieure à cette publication : route, frontière, orientation et
mémoire.

La complétion
\[
 1000=729+243+27+1
\]
montre comment une unité compacte se distribue en espace ambiant ternaire, mot
visible, clôture et unité résiduelle. Elle n'identifie pas à elle seule quatre ères
cosmologiques. Elle fournit une architecture de strates capable de distinguer l'
état profond de sa projection. L'application cosmologique prend la forme
\[
 \mathfrak R_{\rm cos}:
 \text{état HMT de frontière}
 \dashrightarrow
 (M,g,e,\omega,\Psi).
\]
Sa tâche est d'assigner un rang, une échelle, une connexion, un contenu matériel et une loi d'évolution.
La thèse de la distorsion initiale affirme que la grande courbure observée est l'ombre
continue d'une transition de cette application.

La torsion d'Einstein--Cartan fournit un mécanisme dynamique de rebond. Pour une
matière de densité de spin,
\[
 s^2\propto a^{-6}.
\]
La contribution torsionnelle croît vers les petits volumes plus vite que la matière
ordinaire. Une équation effective peut s'écrire
\[
 H^2
 =
 \frac{8\pi G}{3}
 \left(\rho_{\rm m}-\rho_{\rm spin}\right)
 -
 \frac{k}{a^2}
 +
 \frac{\Lambda}{3},
\]
avec les facteurs de \(c\) rétablis selon que l'on emploie des densités de masse ou d'
énergie. L'égalité
\[
 \rho_{\rm m}=\rho_{\rm spin}
\]
localise le candidat d'annulation matérielle--torsionnelle. Il y a rebond lorsque, à une
échelle non nulle, l'expression complète du membre de droite s'annule,
\[
 H^2=0,
\]
et que l'équation d'accélération donne \(\dot H>0\) ; les termes de courbure et
\(\Lambda\) doivent s'annuler, se compenser ou être expressément inclus dans cette condition.
Sous ces hypothèses, la contraction atteint un minimum et s'inverse. Le mécanisme
n'introduit pas de pression arbitraire : il provient de l'équation de Cartan et du
courant de spin.

HMT interprète ce seuil comme clôture et réouverture de feuillet. Pendant la contraction,
l'extension visible se compacte et la mémoire interne augmente. Au point de
retour, l'orientation publiée change et une nouvelle branche expansive devient
accessible. La singularité cesse d'être la seule issue possible parce que la variable d'
état ne se réduit pas au facteur d'échelle. La route enrichie peut continuer lorsque
la carte cosmologique précédente perd du rang.

La même géométrie permet de formuler un univers fils. Pour un observateur extérieur,
une région en effondrement se retrouve derrière un horizon. Pour la section globale, la clôture
peut ouvrir une frontière dotée de sa propre orientation et de son propre temps. Le raccordement est décrit
comme un cobordisme ou une transition entre composantes de feuillet, non comme l'apparition
de matière à partir du néant. L'information qui n'est pas publiée dans l'univers parent
peut devenir une condition de frontière de l'univers fils.

Cette image coordonne trois phénomènes habituellement étudiés séparément : effondrement,
horizon et origine cosmologique. La même opération de clôture qui forme une région
piégée peut produire une réouverture après avoir atteint l'échelle torsionnelle. La même
mémoire qui protège l'information extérieure peut fixer les conditions de la nouvelle
branche. Et le même changement d'orientation qui définit la particule et l'antiparticule à l'échelle
microscopique peut acquérir une réalisation causale macroscopique.

La réalisation complète exige une solution globale. Elle doit contenir un extérieur d'
Einstein, une région de torsion non négligeable, un minimum régulier du facteur d'
échelle, une continuation causale et une application d'états entre frontières. Elle doit respecter
la conservation covariante et produire un observable : distribution des résidus,
corrélations non thermiques, modification du spectre des ondes gravitationnelles ou une
empreinte primordiale. Si l'échelle de rebond contredit la nucléosynthèse, le fond cosmique ou les
limites de torsion, le mécanisme précis est écarté. La thèse du feuillet ne
dispense pas de la dynamique ; elle lui fournit l'objet qui doit évoluer.

\section{Régimes elliptique, parabolique et hyperbolique du TRIT et leur projection sur le secteur sombre}

Le \(\TRIT\) n'oriente pas seulement le passé, le seuil et l'ouverture. Ses trois valeurs déterminent
les trois algèbres quadratiques réelles
\[
 \mathbb A_\tau
 =
 \mathbb R[X]/(X^2+\bar\tau).
\]
Les générateurs satisfont
\[
 \bar\tau=+1:
 \quad J^2=-I,
 \qquad
 \bar\tau=0:
 \quad J^2=0,
 \qquad
 \bar\tau=-1:
 \quad J^2=+I.
\]
Ce sont les régimes elliptique, parabolique et hyperbolique. La géométrie modèle
\[
 g_{s,\rho}
 =
 \dd r^2+S_{s,\rho}(r)^2\dd\theta^2
\]
s'obtient avec
\[
 S_{+,\rho}(r)=\rho^{-1}\sin(\rho r),
 \qquad
 S_{0,\rho}(r)=r,
 \qquad
 S_{-,\rho}(r)=\rho^{-1}\sinh(\rho r),
\]
et possède
\[
 K_G=s\rho^2.
\]
Les trois courbures ne sont pas attribuées simultanément à la même métrique en un même
point. Ce sont des feuillets coordonnés d'un état dont la carte physique sélectionne le régime
réalisé.

L'espace de de Sitter démontre de manière exacte comment une seule géométrie
quadridimensionnelle peut admettre trois présentations spatiales :
\[
 \dd s^2
 =
 -\dd t^2+a_s(t)^2\dd\Sigma_s^2,
 \qquad s\in\{+1,0,-1\}.
\]
Les fonctions
\[
 \begin{aligned}
 a_+(t)&=H^{-1}\cosh(Ht),
 &
 a_0(t)&=\eexp^{Ht},\\
 a_-(t)&=H^{-1}\sinh(Ht),
 &
 t&>0,
 \end{aligned}
\]
satisfont
\[
 \frac{\ddot a_s}{a_s}=H^2,
\qquad
 \left(\frac{\dot a_s}{a_s}\right)^2
 +
 \frac{s}{a_s^2}
 =
 H^2,
\qquad
 R=12H^2.
\]
Un feuilletage fermé, un plat et un ouvert décrivent le même espace-temps.
Planéité spatiale locale, section hyperbolique et accélération positive ne sont pas des signes
incompatibles. Ils peuvent être des feuillets distincts d'une même géométrie globale.

Ce fait fournit un modèle continu pour l'intuition HMT. La valeur publiée
sur un feuillet ne contient pas nécessairement toute la courbure de l'état. Un observateur
qui travaille dans le feuilletage plat peut reconstruire comme source supplémentaire l'effet
provenant de l'incidence de frontière ou d'une autre présentation. La matière sombre et
l'énergie sombre reçoivent ainsi une interprétation unifiée, mais différenciée.

La matière sombre apparente peut correspondre à la réponse gravitationnelle d'états
qui ne possèdent pas de représentant ponctuel sur le feuillet visible. Leur masse se manifeste dans les
orbites et l'effet de lentille parce que la géométrie complète transporte leur incidence, bien que le
lecteur électromagnétique ne publie pas de particule ordinaire. L'énergie sombre
apparente peut correspondre à la réponse globale d'ouverture et de courbure de feuillet,
publiée comme accélération du facteur d'échelle. L'une appartient principalement à la
distribution de rétention non visible ; l'autre, à l'évolution de la frontière
cosmologique.

La charnière
\[
 \frac{G_{\rm tor}}{I_{\rm tor}}=1,
 \qquad
 \frac{G_{\rm av}}{I_{\rm av}}
 =
 \frac{\piHMT}{\phiHMT^2}
\]
indique où chercher ce résidu. Le rapport marqué est une relation structurelle entre
lectures, non un pourcentage cosmologique pouvant être identifié directement à une
densité observée. Pour devenir une théorie du secteur sombre, il doit être intégré
sur la géométrie et produire un tenseur effectif
\(T^{\rm eff}_{\mu\nu}\), ses équations de continuité et une fonction de
transfert.

L'exigence est particulièrement forte parce qu'une seule règle doit traverser plusieurs
échelles. Dans les galaxies, elle doit relier la densité baryonique, la vitesse circulaire et l'effet de lentille
sans introduire un nouveau paramètre pour chaque objet. Dans les amas, elle doit coordonner le gaz,
l'effet de lentille faible et la dynamique. En cosmologie, elle doit prédire \(H(z)\), les distances,
la croissance des structures et les potentiels de lentille. Si la même géométrie multifeuillet
produit ces profils avec des paramètres fixés avant la consultation de chaque système,
HMT--MD aura remplacé deux secteurs inconnus par une seule architecture. Si
elle requiert des ajustements indépendants, l'interprétation n'aura pas atteint son objectif.

Le triple feuilletage impose en outre une discipline conceptuelle. Toute courbure
hyperbolique n'est pas un signe d'énergie sombre, et toute planéité locale n'élimine pas la
courbure globale. La carte, le domaine et l'échelle doivent être déclarés. Le paramètre
\(H\) de de Sitter n'est pas non plus sélectionné par la seule existence des trois
feuilletages. Les mathématiques prouvent leur compatibilité ; le transducteur cosmologique doit
dériver l'échelle physique et son évolution.

\section{Orientation temporelle du fond cosmologique et conservation de la mémoire sous expansion}

La clôture dodécaphasique produit un schéma angulaire très contraint. Soit
\[
 \phi_k=\phi_0+\frac{2\pi k}{12},
 \qquad k=0,\ldots,11,
\]
et
\[
 w_k
 =
 \frac1{12}
 \left[
 1+2\epsilon\cos(3\phi_k-\delta)
 \right],
 \qquad
 |\epsilon|\leq\frac12.
\]
Les poids sont non négatifs et leur somme vaut un. Leur transformée discrète n'a de support que
dans
\[
  m=0,\pm3\quad(\operatorname{mod}\ 12).
\]
Lors de la projection sur les multipôles bas, pour \(\ell=2\), seul peut subsister
\(m=0\) ; pour \(\ell=3\), \(m=0,\pm3\). La famille \(C_{12}\) transforme la roue
dodécaphasique en une enveloppe angulaire exacte.

L'importance physique réside dans son caractère contraignant. L'axe, la phase \(\delta\), l'
amplitude \(\epsilon\), le masque céleste, l'estimateur et l'ensemble des multipôles
doivent être fixés avant d'examiner la carte qui servira à la confrontation. Ensuite, on
compare les données de fréquences et de missions indépendantes à des simulations isotropes
et à des modèles de contamination. Un signal stable dans le support prévu renforcerait
la généalogie dodécaphasique. Une sélection ultérieure de l'axe ou un signal dominant
hors de \(m=0,\pm3\) excluraient cette réalisation précise.

La flèche cosmologique du temps utilise le même lecteur ternaire qui organise les
feuillets :
\[
 +1
 \longmapsto
 \text{retour, mémoire, passé},
\]
\[
 0
 \longmapsto
 \text{seuil, présent},
\]
\[
 -1
 \longmapsto
 \text{ouverture bidirectionnelle, futur}.
\]
Le présent n'est pas un point qui se déplace sur une droite préalablement tracée. C'est
le seuil auquel une ouverture admissible s'actualise et commence à être conservée comme
mémoire. Le passé est le secteur des retours déjà inscrits ; le futur est l'ensemble
des continuations compatibles qui ne se sont pas encore closes.

L'évolution complète peut être réversible tandis que l'histoire publiée possède une
orientation. L'asymétrie apparaît dans la préparation, dans la sélection de frontière et
dans le lecteur qui abandonne une partie de la fibre. L'entropie ne crée pas le temps ; elle mesure l'
information qu'une présentation cesse de distinguer. Pour une action d'histoires
\(I(\gamma)\), la mesure
\[
 \mathbb P_\beta(\gamma)
 =
 \frac{\eexp^{-\beta I(\gamma)}}{Z_\beta}
\]
minimise
\[
 \mathbb E_{\mathbb P}[I]
 -
 \beta^{-1}H(\mathbb P).
\]
À la limite de basse température, l'interaction minimale domine ; à température
finie, la multiplicité des routes intervient. L'action et l'entropie sont deux
composantes d'une même sélection, non deux causes rivales.

Une histoire d'aller et de retour peut être purifiée en
\[
 |\Omega_\beta\rangle
 =
 \sum_\gamma
 \sqrt{\mathbb P_\beta(\gamma)}
 |\gamma\rangle_{\rm ida}
 |\mathcal C\gamma\rangle_{\rm vuelta}.
\]
L'entropie réduite est \(H(\mathbb P_\beta)\). L'univers observé peut posséder
une flèche parce que nous accédons à une branche, tandis que l'état étendu conserve la
conjugaison. Cette image relie l'orientation microscopique de la particule et de l'
antiparticule à l'asymétrie macroscopique sans affirmer que les deux échelles sont
identiques.

La cosmologie HMT--MD est ainsi organisée par une même opération fondamentale :
plier, conserver et déployer. La matière est une section qui retient l'espace ; l'
horizon replie une présentation sur un autre feuillet ; le rebond déploie une branche après
la compactification ; le secteur sombre est la réponse d'une information géométrique que
le feuillet visible ne publie pas ; et le temps est l'orientation d'ouverture et de retour.
Chacun de ces énoncés acquiert une force physique en partageant le même domaine d'
état, non en formant une collection d'analogies.

La théorie s'engage sur des conséquences claires. L'extérieur du trou noir doit
suivre la relativité dans son régime éprouvé et l'inversion doit apporter une
continuation et un signal. Le rebond doit fixer son échelle par la densité de
spin et respecter la nucléosynthèse et le fond cosmique. La fonction de transfert
multifeuillet doit expliquer conjointement la dynamique, l'effet de lentille et l'expansion. Le schéma
\(C_{12}\) doit résister à une confrontation prospective. La flèche temporelle doit
respecter la réversibilité globale et le coût d'effacement. Ces conditions ne diminuent pas la
thèse cosmologique : elles sont le chemin par lequel une architecture mathématique devient
une physique de l'univers.

